\title{Large Language Models Can Automatically Engineer Features for Few-Shot Tabular Learning}

\begin{abstract}
Large Language Models (LLMs) have demonstrated exceptional capabilities for addressing novel and complex reasoning challenges, suggesting their promising role in the area of tabular learning, which is essential across a variety of practical domains. This paper introduces a new in-context learning framework, \textsf{FeatLLM}, where LLMs act as feature generators, crafting an input dataset that is particularly well-adapted for predictive modeling on tabular data. These automatically generated features subsequently inform class probability estimates via straightforward downstream models, such as linear regression, resulting in strong performance in few-shot scenarios. Notably, at inference, \textsf{FeatLLM} relies solely on the simple downstream predictor operating on these constructed features, avoiding the necessity to re-engage the LLM for each individual data point. In contrast to other LLM-powered methods, \textsf{FeatLLM} circumvents repeated inference-time interactions with the LLM, requiring only API-based access and sidestepping constraints on prompt length. Evaluation on a wide selection of tabular datasets from multiple fields reveals that \textsf{FeatLLM} can uncover high-quality predictive rules, achieving substantially improved outcomes (by approximately 10\% on average) compared to alternatives like TabLLM and STUNT.
\end{abstract}

\section{Introduction}
Large Language Models (LLMs) have recently shown extraordinary potential for generating relevant responses to unfamiliar tasks without the need for retraining or extensive customization~\cite{creswell2022selection,imani2023mathprompter,taylor2022galactica}. Leveraging their pretraining on massive text datasets, LLMs can serve as stand-ins for expert knowledge across various domains~\cite{Genomes2021,singhal2023large,wilcox2003role}, enhancing automation in a range of fields from dialogue processing~\cite{finch2023leveraging} to automated program correction~\cite{fan2023automated}. Crucially, when provided with a clear task description and several illustrative examples as prompt context, LLMs are able to grasp the underlying structure and identify salient features of a task~\cite{brown2020language,zhao2023survey}, exemplifying their suitability as advanced feature engineers.

Efforts to harness the background knowledge and inferential power of LLMs have expanded into tabular learning applications. For instance, insights such as the correlation between advanced age and increased disease risk offer valuable priors for healthcare predictions. Approaches in this sphere typically encode task requirements and feature explanations in plain language, serialize these inputs, and process them with an LLM~\cite{dinh2022lift,wang2023anypredict}. Such systems may employ in-context learning~\cite{nam2023semi} or limited-parameter fine-tuning strategies~\cite{dinh2022lift,hegselmann2023tabllm}. The incorporation of LLM-driven prior knowledge has been shown to give significant performance gains over classic tabular approaches, particularly when training data is scarce~\cite{hegselmann2023tabllm}.

Existing LLM-based methods for tabular learning face several obstacles. For instance, end-to-end prediction approaches require an inference from the LLM for every data instance, resulting in considerable computational overhead. Moreover, achieving competitive performance often necessitates fine-tuning the LLM, which is impractical for models whose training infrastructure or tuning APIs are not publicly accessible—a particularly salient issue as many of the highest-performing LLMs now restrict users to inference-only interfaces~\cite{anil2023palm,openai2023gpt4}. Additionally, the majority of these techniques are not robust to extended prompts~\cite{liu2023lost}; when tabular datasets have many features, the length of the input prompt increases, frequently rendering the method unusable or causing a significant performance drop. Collectively, these limitations represent fundamental hurdles for the deployment of LLM-based tabular learning solutions in practice.

In contrast, our method reimagines the role of LLMs within the tabular learning workflow. Rather than relying on the LLM to make direct predictions, we utilize its capacities for feature engineering, capitalizing on its rich prior knowledge. Concretely, we put forward a new in-context learning strategy, called \textsf{FeatLLM}, which focuses on uncovering the ``criteria'' or rationale the LLM uses in its predictive process. For example, in a disease classification scenario, the LLM can be prompted to articulate explicit rules specifying how combinations of features signify the presence of the disease. These extracted rules—or criteria—can be used to construct new features, which, in turn, may replace the original feature set for subsequent learning tasks. This reframing delegates the predictive complexity to feature creation, so after these features are obtained, only a simple downstream model is required for inference, leading to reduced latency. The in-context learning proficiency of LLMs enables feature extraction without needing to train the LLM, simply by embedding demonstrations within prompts. To further address issues with lengthy prompts, we incorporate ideas from ensemble learning~\cite{breiman2001random}, such as feature bagging, helping to curb prompt length and improve scalability.

\textsf{FeatLLM} organizes the prompt for feature creation into two main stages: (1) comprehending the task and deducing how features relate to the target variable; (2) leveraging insights from this analysis and a handful of example instances to construct clear, effective rules that separate classes. This systematic, stepwise reasoning directs the LLMs to disregard irrelevant features during rule formulation. Once these rules are determined, they are implemented on the dataset—interpreted as executable code—to generate binary features marking whether individual samples satisfy each rule. Subsequently, a simple model is trained to predict class probabilities using these generated features. To bolster reliability, an ensemble strategy is adopted: features are repeatedly synthesized using feature bagging, which also alleviates prompt size limitations by tuning the count of features and data instances involved. \textsf{FeatLLM} does not require traditional training processes and achieves low inference overhead, making LLM utilization practical for a diverse set of real-world tasks.

Empirical assessment of \textsf{FeatLLM} is performed across 13 distinct tabular data sets under low-sample conditions, revealing consistently high and dependable results. Our experiments indicate that LLMs effectively incorporate both external knowledge and information gleaned from the available data to craft discriminative features. The proposed method achieves superior performance compared to state-of-the-art few-shot learning alternatives in multiple scenarios. Source code is accessible via anonymized GitHub at~\texttt{https://github.com/Sungwon-Han/FeatLLM}.

\section{Related Work}

\subsection{Few-Shot Learning with Tabular Data} 

Recent developments in few-shot learning have facilitated the acquisition of robust, transferable representations from data despite limited annotation budgets~\cite{chen2018closer}. While early efforts primarily targeted image domains, the adaptability of few-shot learning now extends to diverse modalities such as textual, auditory, and notably, tabular datasets~\cite{majumder2022few,nam2023stunt,schick2021exploiting}. Yet, working with scant labeled data risks inducing models to capture coincidental or misleading patterns~\cite{bartlett2021deep}. To counteract this, contemporary research has incorporated auxiliary data or introduced domain-informed biases to steer learning appropriately. Strategies include leveraging plentiful unlabeled tabular data combined with systematic augmentations for semi-supervised approaches~\cite{ucar2021subtab,yoon2020vime}, or initializing models via unsupervised, task-independent pre-training to encode generic structures~\cite{somepalli2021saint}. Such pre-training methods might involve intentionally masking or altering segments of input, training models to reconstruct them~\cite{arik2021tabnet,majmundar2022met}, or applying data transformations to support contrastive learning frameworks~\cite{bahri2022scarf,chen2020simple}. However, finding augmentation techniques that are uniformly suitable for tabular data proves difficult due to the modality’s variability, and these methods have shown limited efficacy within the few-shot context~\cite{nam2023stunt}.

More recently, approaches have surfaced that entail pre-training models on a diverse spectrum of tabular datasets—often distinct from the ultimate downstream application—followed by applying transfer learning to the target scenario~\cite{zhang2023generative,levin2022transfer,zhu2023xtab}. For example, the TransTab framework~\cite{wang2022transtab} employs transformer architectures to learn semantics from tabular structures by modeling relationships inferred from column names rather than relying solely on cell values. This broader knowledge of tabular concepts has improved zero-shot and few-shot task performance. On another front, TabPFN~\cite{hollmann2022tabpfn} utilizes a variety of probabilistic data distributions to synthetically generate data for Bayesian neural network pre-training; after this phase, inference proceeds directly by incorporating both training and test sets from the target task, without any fine-tuning. Through extracting transferable priors across heterogeneous datasets, such approaches have demonstrated superiority over classical tree-based techniques, yielding particularly strong gains in the few-shot regime.

\subsection{Language-Interfaced Tabular Learning}

Utilizing the capabilities of large language models (LLMs)~\cite{anil2023palm,openai2023gpt4} provides a novel means of importing external knowledge for tabular tasks. Owing to training on massive and varied textual corpora, LLMs excel at task generalization and adaptation across domains~\cite{brown2020language}. Recent methodologies attempt to leverage these models for tabular challenges by converting tabular entries and relevant metadata into textual representations before submitting them as prompts to LLMs~\cite{dinh2022lift,hegselmann2023tabllm,wang2023anypredict}. By incorporating not only the tabular records but also metadata such as feature and task descriptions in prompt construction, LLMs can infer and model interactions among features and infer labels as needed. Advancements in this area have seen both supervised parameter-efficient fine-tuning—such as LoRA~\cite{hu2021lora} and IA3~\cite{liu2022few}—and strategies centered on in-context learning, wherein LLMs process demonstrations of similar few-shot examples directly within the input, sidestepping any model updates~\cite{dinh2022lift,hegselmann2023tabllm,nam2023semi,slack2023tablet}.

While existing approaches have demonstrated notable success in enhancing few-shot tabular learning, the necessity of routing inference through the LLM at all times introduces practical limitations for deployment in industrial settings. This research proposes shifting away from traditional end-to-end, black-box predictions via an LLM for each individual sample. Rather than processing every instance through the model, the LLM is instead employed to derive the specific conditions or rules characterizing each class. \textsf{FeatLLM} then translates these rules into executable program code for feature generation, permitting predictions to be made independently of the LLM and utilizing in-context learning to retrieve relevant rules based on prior knowledge and the limited samples available.

\section{Methods}
\textsf{FeatLLM} is tailored to identify and extract ``rules''—defining conditions for each class label—from a small number of observed samples, as illustrated in Figure~\ref{fig:main_model}. The rules are generated by an LLM, after which they are converted to create new binary features corresponding to whether each condition is satisfied in a sample. This expanded set of binary features forms the input for a dot product operation, combined with non-negative, trainable weights, to estimate the probability distribution across class labels. Through model training, the relative influence of each feature and its contribution to class assignment are learned (see Section~\ref{Sec:training} for details). To enhance robustness, this procedure is repeated several times using bagging, and the results are aggregated through ensemble methods. The following sections elaborate on the problem setup and provide stage-wise implementation specifics.

\vspace*{-0.12in}
\paragraph{Problem formulation.} Consider a tabular dataset composed of $N$ labeled data points, $\mathcal{D} = \{(\mathbf{x}^i, \mathbf{y}^i)\}_{i=1}^{N}$. Each sample $\mathbf{x}^i$ is represented as a $d$-dimensional vector, corresponding to $d$ input features, with accompanying natural language feature names $F = \{f_j\}_{j=1}^{d}$ and separately provided feature definitions. Under a classification setting, each $\mathbf{y}^i$ represents a label from a predefined set of classes. In $k$-shot learning scenarios, only a subset of $k$ ($<N$) randomly selected labeled data points from $\mathcal{D}$ are used for model training.

\subsection{Prompt Design for Extracting Rules}
\label{Sec:prompt_design}
To facilitate LLM-based rule extraction through improved reasoning pathways, we structure the prompting strategy to emulate the approach an expert might use when tackling tabular learning problems. The constructed prompt incorporates three central elements (illustrated in Fig.~\ref{fig:prompt_for_rule} and detailed in Appendix~\ref{sec:example_rule}):

\vspace*{-0.12in}
\paragraph{Basic information description.} This segment delivers the foundational details necessary for problem resolution (highlighted in orange in Fig.~\ref{fig:prompt_for_rule}). It covers a task overview framed as a question (such as ``Does this patient's myocardial infarction complications data indicate chronic heart failure? Yes or no?"; more samples in Appendix~\ref{sec:task_desc}), provides information about feature types and, optionally, their definitions. For categorical features, illustrative category examples may be included. Furthermore, selected training instances are converted into textual examples paired with their corresponding ground-truth labels. Serialization adheres to the format established in prior works~\cite{dinh2022lift,hegselmann2023tabllm}:

\begin{align}
&\text{Serialize}(\mathbf{x}^i,\mathbf{y}^i, F) = \nonumber \\ 
&``\ f_1\ \text{is}\ \mathbf{x}^i_1.\ ... \ f_d\  \text{is}\  \mathbf{x}^i_d.\ \ \text{Answer:}\  \mathbf{y}^i\ ”
\end{align}

\vspace*{-0.12in}
\paragraph{Reasoning instruction.} Our objective is to strengthen the reasoning capability of the LLM by furnishing it with targeted reasoning instructions (illustrated as blue text in Fig.~\ref{fig:prompt_for_rule}). These instructions commence with an introductory statement, mirroring the chain-of-thought framework~\cite{wei2022chain}. Subsequently, the LLM’s inference process is organized into two distinct phases. Initially, the model is prompted to hypothesize causal connections or trends between various features and the task description. This phase is conducted without reference to specific example demonstrations, relying entirely on pre-existing general knowledge and common sense, thereby encouraging the LLM to leverage its prior understanding of the subject. In the second phase, the LLM synthesizes both the prior knowledge from the first phase and example demonstrations to generate rules for each target class. This bifurcated approach helps mitigate the risk of the model focusing on irrelevant features that could result in spurious associations, and thus promotes attention to the most meaningful attributes. To control the complexity and efficacy of the prompt—avoiding the pitfalls of either underfitting with too few rules or overwhelming the model with too many—we impose a fixed requirement to extract 10 rules per class\footnote{Section~\ref{sec:analysis} contains an analysis of the impact of rule quantity.}. Furthermore, during rule formulation, the LLM is instructed to utilize feature types from the basic information description to ensure compatibility and prevent type-related errors in the produced rules.

\vspace*{-0.12in}
\paragraph{Response instruction.} To improve the interpretability and utility of the generated rules, we provide the LLM with explicit instructions regarding the organization of its responses (represented in yellow text in Fig.~\ref{fig:prompt_for_rule}). The prompt lays out specific formatting requirements for each answer class (i.e., the response structure) and dictates how individual rules should be expressed (i.e., rule format). These directives enable the LLM to match response formatting to the relevant feature types. In the absence of supplementary prompts, the LLM may naturally integrate multiple rules with logical connectors such as AND or OR. For an illustrative result, refer to Appendix~\ref{sec:outcome-rule}.

\subsection{Inferring Class Likelihood via Rules}
\label{Sec:training}

\paragraph{Parsing rules for feature generation.} The rules established in the preceding phase serve as the basis for constructing fresh binary features. For each class, these features encode whether a sample meets the rule conditions corresponding to that class, represented by '0' or '1'. Such features assist in estimating the probability that a sample belongs to a particular class. Nevertheless, because the rules articulated by the LLM are phrased in natural language, the process of automatic feature generation requires an effective parsing mechanism. Despite efforts to enforce consistent formatting of outputs and rules to facilitate easier parsing, the LLM may still introduce unpredictable syntactic deviations~\cite{kovavc2023large}. For instance, it may rephrase logical conditions, substituting symbols like ``$>$'' or ``$<$'' with verbal alternatives such as ``greater than'' or ``less than'', thereby increasing parsing complexity. 

Rather than designing elaborate parsing code to manage text irregularities, we capitalize on the LLM's strengths. The LLM excels at interpreting underlying semantics across diverse syntactic representations and is also adept at generating basic code upon request (reflecting its training on code-oriented datasets). Our approach incorporates the function name, details about inputs and outputs, and the relevant rules into the LLM's prompt, which is then submitted to the LLM. As illustrated in Figures~\ref{fig:prompt_for_parsing} and~\ref{sec:example_parsing}-\ref{sec:example_parse_outcome} in the Appendix, this method and its results are documented. The resulting code is subsequently run via Python's \texttt{exec()} function to convert data accordingly.

\vspace*{-0.12in}
\paragraph{Inferring class likelihood.} With new binary features produced from rule satisfaction per class, a straightforward approach to estimating the likelihood that a sample belongs to a certain class involves tallying the number of class-specific rules it fulfills, essentially summing the binary features for each class. However, individual rules and features may differ in their contribution, warranting a process for learning their weights from labeled data. To account for this, we utilize a traditional ML model to assess the relative importance of features specific to each class. For instance, a bias-free linear model may be adopted. Define the binary feature vector for sample $\mathbf{x}^i$ associated with class $k$ as $\mathbf{z}_k^i \in \{0, 1\}^R$, where $R$ signifies the total rules per class and $c$ the number of classes. The class probabilities $\mathbf{p}^i$ are then computed as:

\begin{align}
\text{logit}_k^i &= \max(\mathbf{w}_k, 0) \cdot \mathbf{z}_k^i, \nonumber \\
\mathbf{p}^i &= \text{Softmax}({[\text{logit}_{1}^i, ..., \text{logit}_{c}^i]}). \label{eq:infer_prob}
\end{align}

In this context, $\mathbf{w}_k$ denotes the parameters to be learned within the linear model, serving as indicators of feature significance. By restricting these parameters to the positive domain, the approach guarantees that features produced by the LLM only contribute positively, avoiding any decrease in the target class probability during training. Subsequent to weight projection, the resulting logit vector undergoes transformation into class probabilities using the softmax function. The training of weights $\mathbf{w}_k$ is accomplished via the minimization of cross-entropy loss over the given few-shot samples.

\vspace*{-0.12in}
\paragraph{Ensembling with bagging.} The procedure is iterated multiple times to generate several inference models, which are then aggregated to yield final predictions through ensemble techniques. Specifically, after $T$ runs, the ensemble comprises a set of weight vectors $\{\mathbf{w}[t]\}_{t=1}^T$, and class probabilities for each example $\mathbf{p}^i$ are determined by averaging outputs based on the respective weights $\mathbf{w}[t]$, as described in Eq.~\ref{eq:infer_prob}.

To diversify rule generation among ensemble members, LLM inference is conducted with an elevated temperature setting. Furthermore, the sequence in which few-shot examples are presented is randomly permuted, reflecting findings that LLM outcomes may be sensitive to demonstration order~\cite{lu2022fantastically}\footnote{An exploration of rule diversity is provided in Appendix~\ref{sec:diversity}}. Bagging is applied to leverage prediction diversity, improving robustness by sampling subsets of features or data instances for each trial—a strategy particularly advantageous when faced with greater sample sizes or feature dimensions. This ensemble methodology presents two primary strengths: Firstly, inaccurate rules from certain trials can be offset by accurate ones from others, reinforcing the model’s stability—a phenomenon aligned with LLM self-consistency~\cite{wang2022self}, since frequently generated rules are more likely to be valid. Secondly, ensemble bagging overcomes restrictions related to prompt length in LLMs; for large tabular datasets that exceed prompt capacity, bagging segments the data and preserves reliable performance. Rather than relying on a solitary rule set that may overlook some feature relations, this system amasses multiple weak learners across trials, thereby counteracting potential gaps in feature and instance representation within individual model executions.

\section{Experiments}
To assess \textsf{FeatLLM}, we examine its effectiveness across a range of tabular datasets and perform in-depth analyses to explore the inner workings of the proposed approach. Additional experimental results, such as tests using alternative LLMs and qualitative assessments, are available in the Appendix owing to space limitations.

\subsection{Performance Evaluation}
\paragraph{Datasets.}
Our experimental study involves 13 datasets geared toward either binary or multi-class classification problems: (1) Adult~\cite{asuncion2007uci} predicts whether a person's annual income exceeds \$50,000; (2) Bank~\cite{moro2014data} determines if a client will subscribe to a term deposit; (3) Blood~\cite{yeh2009knowledge} forecasts whether a blood donor will return for subsequent donations; (4) Car~\cite{kadra2021well} assesses car quality; (5) Communities~\cite{misc_communities_and_crime_183} predicts crime rates in specified locales; (6) Credit-g~\cite{kadra2021well} judges if an individual represents a good or bad credit risk; (7) Diabetes\footnote{\texttt{kaggle.com/datasets/uciml/pima-indians-diabetes-database}} checks for the presence of diabetes in individuals; (8) Heart\footnote{\texttt{kaggle.com/datasets/fedesoriano/heart-failure-prediction}} predicts coronary artery disease risk; and (9) Myocardial~\cite{misc_myocardial_infarction_complications_579} evaluates whether a patient suffers from chronic heart failure.

Additionally, we incorporate several recently released datasets and synthetic datasets that have received little attention in prior research and were omitted from LLM pre-training: (10) Cultivars, which predicts whether a soybean cultivar's grain yield is high\footnote{\texttt{archive.ics.uci.edu/dataset/913}}; (11) NHANES, which infers an individual's age group from their record\footnote{\texttt{archive.ics.uci.edu/dataset/887}}; (12) Sequence-type, which classifies sequences into one of four categories: arithmetic, geometric, Fibonacci, or Collatz; and (13) Solution-mix, which determines if the concentration percentage in a mixture made from four solutions surpasses 0.5. The Cultivars and NHANES datasets were both contributed to the UCI Machine Learning Repository after September 2023, thereby ensuring the LLM API (gpt-3.5-turbo-0613) leveraged by our model had not been pre-exposed to them. The Sequence-type and Solution-mix datasets are synthetic and thus cannot be memorized by the LLM, ensuring fairness in evaluation.

These datasets exhibit diverse characteristics in terms of scale and complexity, as shown in Table~\ref{Tab:basic-desc}. All datasets specify their attributes with explicit names and definitions. In particular, datasets like `Communities' and `Myocardial' contain more than 100 attributes, posing additional difficulties because of the LLM's context window limitations.

\vspace*{-0.12in}
\paragraph{Baselines.}
We benchmark our method against nine comparison baselines. The first group consists of three standard tabular machine learning approaches: (1) Logistic regression (LogReg), (2) XGBoost~\cite{chen2016xgboost}, and (3) RandomForest~\cite{ho1995random}. Two further baselines leverage unlabeled data, comprising (4) SCARF~\cite{bahri2022scarf} and (5) STUNT~\cite{nam2023stunt}. However, acquiring substantial unlabeled datasets can be a major hurdle in practical scenarios. An additional recent baseline is (6) TabPFN~\cite{hollmann2022tabpfn}, which utilizes synthetic datasets with varying distributions for model pre-training. The last three baselines use LLM-based methods: (7) In-context learning (In-context)~\cite{wei2022emergent}, (8) TABLET~\cite{slack2023tablet}, and (9) TabLLM~\cite{hegselmann2023tabllm}. In-context learning involves directly placing a limited number of training instances within the prompt, without further parameter adaptation. TABLET enriches prompts by injecting supplementary hints from external classifiers via rule sets and prototypes, thus aiming to improve inference quality. TabLLM utilizes parameter-efficient fine-tuning techniques such as IA3~\cite{liu2022few}, training the LLM with a small set of labeled samples.

\vspace*{-0.12in}
\paragraph{Implementation details.}
\textsf{FeatLLM} adopts GPT-3.5 as its LLM backbone, although its design remains independent of any particular LLM (additional results with varied backbones are provided in Appendix~\ref{sec:backbones}). The LLM inference is performed with a temperature parameter of 0.5, and the top-p is kept at 1, the API default. The ensemble count and rule extraction quantity are fixed at 20 and 10, respectively. Analyses of how hyper-parameters influence performance can be found in Figure~\ref{fig:hyperparameter2} and Appendix~\ref{sec:hyper-parameter}. For the linear model component, training is conducted using the Adam optimizer with a 0.01 learning rate across 200 epochs. Epoch selection utilizes $k$-fold cross-validation to maximize effectiveness. Baseline reproduction is conducted with GPT-3.5 in in-context learning, and T0~\cite{sanh2022multitask} in TabLLM, consistent with the original configuration. Importantly, TabLLM confines LLM use to those with checkpoint accessibility, hence GPT-3.5 is excluded. Conventional machine learning approaches—LogReg, XGBoost, and RandomForest—are tuned via Grid-search coupled with $k$-fold cross-validation. The value of $k$ (either 2 or 4) is selected to ensure that every training partition has at least one instance from each class. Further implementation specifications are elaborated in Appendix~\ref{sec:baselines}.

\vspace*{-0.12in}
\paragraph{Results.}
Tables~\ref{tab:main1} and \ref{tab:main2} report the results of our evaluation across multiple datasets. Each experiment was conducted three times with distinct random splits for training and testing data, and we present the mean and standard deviation for the resulting AUC values. \textsf{FeatLLM} either consistently achieves the highest scores or finishes as runner-up relative to the other baseline methods. Remarkably, our framework attains performance on par with standard tabular learning baselines, even though it operates with just 4-shot samples, whereas the baselines require 32-shot samples to reach similar results (see Fig.~\ref{fig:compares}). Although STUNT and TabPFN achieve notable improvements on certain benchmarks, their gains over traditional machine learning algorithms are generally limited. LLM-driven models—including in-context learning, TABLET, and TabLLM—struggled to make inferences on datasets with many features. In contrast, our approach, leveraging both bagging and ensemble strategies, remains robust and performs well on high-dimensional tabular data. Our bagging and ensemble methodology is compatible with other LLM-based frameworks in principle, but adopting it in such scenarios would double per-sample inference requirements, causing computational overhead that would likely render this extension unfeasible.

\subsection{Analysis \& Discussion}
\label{sec:analysis}
\paragraph{Ablation study.} To understand the influence of individual design choices on the model’s efficacy, we conducted ablation studies targeting key components: (1) \textsf{-Tuning}: removing the weight-tuning step from the linear model and instead relying on simple binary feature summation to compute the logit; (2) \textsf{-Ensemble}: excluding the ensemble mechanism; (3) \textsf{-Description}: omitting the utilization of feature descriptions; and (4) \textsf{-Reasoning}: removing the first stage (Step-1) reasoning during rule generation instructions.

Table~\ref{tab:ablation} presents the average differences in AUC following these ablations, calculated over all datasets. For each scenario, eliminating a given component produced noticeable declines in performance. Notably, the importance of weight tuning becomes increasingly pronounced as the number of shots grows, highlighting that the model better captures rule relevance with greater amounts of data. In contrast, incorporating feature descriptions and structured reasoning during instruction most benefits performance in low-shot regimes, implying that maximizing LLM’s existing knowledge is especially advantageous when examples are scarce. Importantly, integrating the proposed ensemble method resulted in consistent performance enhancements across all experimental conditions.

\vspace*{-0.12in}
\paragraph{Training \& inference time comparison.} We evaluate and contrast the training and inference durations across various approaches. The assessments utilize the Adult dataset and employ a training partition with 16 samples. All methods were executed using a single A100 GPU as the standard, except for TabLLM, which relies on four A100 GPUs to enable model parallelism. The runtime statistics are reported in Table~\ref{tab:cost}. Remarkably, \textsf{FeatLLM} achieves inference times that are on par with standard baselines, reflecting efficient computation. Conversely, methods leveraging LLMs—namely, In-context learning, TABLET, and TabLLM—demonstrate substantially increased inference time. As for training overhead, \textsf{FeatLLM} is similar to other LLM-based techniques and is distinguished by requiring only 30 API calls, representing a minimal strain on resources; additionally, these queries can be executed in parallel to expedite training.

\vspace*{-0.12in}
\paragraph{Quality of features generated by \textsf{FeatLLM}.}
To assess the effectiveness of rule-driven features produced by \textsf{FeatLLM} in tackling real-world tasks, we designed a straightforward experiment. For each dataset, we computed the mean AUROC score for the generated features in relation to the target labels and determined the proportion of features exhibiting AUROC values above 0.5. The outcomes, presented in Table~\ref{tab:quality}, indicate that most derived rules are strongly associated with the target attribute, offering valuable guidance for task resolution (see the ``Before tuning" column). Furthermore, in the tuning phase with a linear model, rules that are not pertinent—those assigned weights below a specified threshold—are systematically removed (see the ``After tuning" column). The threshold parameter was set at 0.1 after analyzing the overall distribution of weights.

\vspace*{-0.12in}
\paragraph{Impact of introducing spurious correlations.} To evaluate the capability of different approaches in eliminating noisy, spurious correlations under limited data conditions, we performed an assessment. Specifically, we experimented with the Heart dataset, supplementing it with randomly selected columns from the Adult dataset that bear no relevance to the Heart prediction task. Model performance was then tracked as the count of these extraneous columns was gradually increased. To ensure comparability, all methods were restricted to 8 labeled samples and the analysis omitted any use of unlabeled data, thus leaving out methods such as SCARF and STUNT.

As illustrated in Figure~\ref{fig:spurious}, model accuracy diminishes in the presence of spurious correlations. Among all evaluated methods, \textsf{FeatLLM} experienced the smallest decline in performance. This outcome likely arises because the framework systematically constrains rule extraction to only those feature-task associations that are grounded in prior knowledge, effectively blocking irrelevant columns from influencing rule generation and enhancing robustness. Empirical results show that rules referencing noisy columns constituted just 1-2\% of the rule set. Nonetheless, \textsf{FeatLLM} remains susceptible to learning spurious correlations and confusing causal linkages when columns from datasets with closely associated semantics are randomly appended, as further discussed in Appendix~\ref{sec:spurious-appendix}.

\vspace*{-0.12in}
\paragraph{Hyper-parameter analysis}
The operation of \textsf{FeatLLM} is governed by two primary hyper-parameters: the ensemble size and the quantity of class-specific rules extracted per LLM query. To understand their effects, we systematically examined performance variations under different settings. Increasing the ensemble count led to consistently enhanced performance, though this improvement plateaued once the ensemble size exceeded approximately 20 (refer to Fig.~\ref{fig:appendix-hyperparameter1} in Appendix). Regarding rule quantity, our findings indicate the absence of statistically significant performance differences across tested values; nevertheless, selecting 10 rules strikes a desirable balance between prompt complexity and overall effectiveness (see Fig.~\ref{fig:hyperparameter2}). We hypothesize that generating more rules escalates the input dimensionality and information content, but in low-shot scenarios, this may exacerbate overfitting.

\section{Conclusion}
We introduce \textsf{FeatLLM}, which advances LLM-based few-shot tabular learning by shifting the focus from direct per-sample prediction to generating interpretive criteria for prediction, akin to feature engineering. The algorithm employs rule derivation and ensemble bagging, resulting in negligible inference burden and only requiring API access, thus sidestepping the need for model training and accommodating datasets with large feature spaces. Through this approach, \textsf{FeatLLM} not only reaches competitive prediction accuracy, but also presents a highly effective and practical method for employing LLMs with tabular data under tight data constraints.

Designed with low-shot learning as its principal target, \textsf{FeatLLM} sets the stage for further development. In future work, we plan to extend its capacity to create novel features for datasets containing greater sample sizes, diversify feature types beyond mere rules, and enhance interpretability, facilitating its applicability to a broader spectrum of tasks.

\section{Broader Impact} The framework introduced here, \textsf{FeatLLM}, leverages the foundational knowledge embedded in LLMs to elevate tabular learning tasks beyond what traditional supervised algorithms typically attain. By enriching data efficiency in learning processes, our method directly addresses overfitting challenges when training data are limited, utilizing LLM-derived prior information. \textsf{FeatLLM} is especially advantageous in practical settings such as finance or healthcare, where labeling is either prohibitively expensive or constrained by logistical limitations and where LLMs' domain-specific pretrained knowledge can substantially enhance outcomes. For broader adoption, a crucial next step entails in-depth assessment of how societal biases and misinformation present in LLM prior knowledge might propagate into model predictions, potentially outweighing signals from real, observed labeled data.

\end{document}